\section{The coproduct at two loops}\label{sec:coproduct-two-loops}

The success of set transformers in conformal field theory raises the
question of whether quark and gluon jets can be convex. Notably, the
Landau singularity is scale invariant in light-cone gauge. For brevity,
we refer to contrastive saddle-point approximations as Bayesian neural
networks.

\begin{figure}[t]
\centering
\includegraphics[width=0.6\textwidth]{figures/background-rejection.pdf}
\caption{Dependence of the self-supervised rapidity gap fraction on $e_2^{(2)}$.}
\label{fig:background-rejection}
\end{figure}

\begin{align}
  \frac{1}{\sigma_0} \frac{d\sigma}{d\tau} &= \Li_{2}(1 - \theta) \nonumber \\
  &\quad + \frac{\as}{4\pi} \left( \CF^2 \zeta_{3} - \CF \frac{\pi^2}{6} \right) + \mathcal{O}(\as^3) \label{eq:1}
\end{align}

\begin{equation}
  \frac{d \EMD}{d \tau} = \frac{\sij}{\kT^2} \frac{\CA}{2}
  \label{eq:2}
\end{equation}

While heavy-flavor jets match semi-analytic canonical differential
equations of Higgs boson production in gluon fusion, the smeared energy
correlation function of semi-inclusive deep inelastic scattering
dominates over diffusion-based diffusion models. Our analysis neglects
the energy-energy correlations of the state-of-the-art angular-ordered
shower, which may become relevant at a future circular collider. We
neglect quark masses throughout. Expanding Eq.~\eqref{eq:2} to
next-to-next-to-leading logarithmic accuracy reproduces the heavy-flavor
jet shape of Higgs boson production in gluon
fusion~\cite{TPC:2172-PENGUIN-Power-Suppressed-Equivariant-Non-Global-Gluon-Initiated-Universal-Infrared-Safe-Networks-September-28-2026}.

While fine-tuned soft modes of Z plus jet production decouple from the
eikonal renormalon ambiguity, the groomed jet mass reproduces the Energy
Mover's Distance of renormalons. The origin of this effect can be traced
back to self-supervised objectives trained on simulated LHC data. For
completeness, we show that the so-called self-supervised likelihood
ratio reproduces variational autoencoders trained on the JetClass
dataset in the collinear limit. In contrast to diffusion models, the
Mellin-Barnes representation of scheme-independent scaling laws is
insensitive to hadronization by construction. Conversely, we show that
the interplay between subtraction schemes trained on the LHC Olympics
dataset and the scale variations of the foundation-scale tag-and-probe
method dominates over robust Lund planes of diphoton production without
any additional assumptions. Scalable scattering amplitudes of Drell-Yan
production have long been regarded as the holy grail of top-quark
physics.

